%
% session-model.tex
%
% Z Specification for Biff Session Identity Model
%
% Models the full organizational hierarchy: orgs, repos, machines,
% clones, users, sessions, and processes (human operators and AI agents).
%
% Type-check: fuzz session-model.tex
%

\documentclass[a4paper,10pt,fleqn]{article}
\usepackage[margin=1in]{geometry}
\usepackage{fuzz}
\usepackage[colorlinks=true,linkcolor=blue,citecolor=blue,urlcolor=blue]{hyperref}

\begin{document}

\title{Biff Session Identity Model:\\A Z Specification}
\author{Punt Labs}
\date{March 2026}
\hypersetup{
  pdftitle={Biff Session Identity Model: A Z Specification},
  pdfauthor={Punt Labs},
  pdfsubject={Z Specification},
  pdfcreator={fuzz}
}
\maketitle

\tableofcontents
\newpage

%%--------------------------------------------------------------------
\section{Introduction}
%%--------------------------------------------------------------------

This specification models the organizational hierarchy for biff's
communication system.  The hierarchy spans from organizations down to
individual processes within terminal sessions.

\medskip
\noindent\textbf{Containment hierarchy:}

\begin{enumerate}
  \item \textbf{Organization} --- owns repositories and has member users.
  \item \textbf{Machine} --- a host that holds clones of repositories.
  \item \textbf{Clone} --- a copy of a repository on a machine.
  \item \textbf{Session} --- a terminal on a machine, working on a branch
    of a cloned repository.  May have zero or one human operator.
  \item \textbf{Process} --- a human operator or AI agent running within
    a session.  Processes form a tree: a main agent spawns subagents.
\end{enumerate}

\noindent\textbf{Design principles:}

\begin{itemize}
  \item The session is the container; humans and agents are both processes
    within it.
  \item A session may have zero or one human operator.  Headless sessions
    (CI, SDK, agent teams) have none.
  \item A session may contain many agent processes, forming a tree rooted
    at top-level agents.
  \item A user may have sessions on multiple machines, branches, and
    worktrees simultaneously.
\end{itemize}

%%--------------------------------------------------------------------
\section{Types}
%%--------------------------------------------------------------------

\begin{zed}
  [NAME, HOST, SESSION, PROCESS]
\end{zed}

$NAME$ identifies organizations, repositories, users, and branches.
These are all string-valued in the implementation and share a namespace.
$HOST$, $SESSION$, and $PROCESS$ are opaque identifiers for their
respective entity types: machine hostnames, terminal session IDs, and
process IDs.

\medskip\noindent
$SESSION$ is the Claude Code \texttt{session\_id}.  It is opaque, it is
durable, and --- the fact the routing layer below is built on --- it is
\emph{stable across} \texttt{claude --resume} and \texttt{--continue},
and \emph{fresh on} \texttt{--fork-session}.  A resumed session recovers
the same \texttt{session\_id} (the recovered identifier is reused, not
minted afresh); a forked session is a genuinely new node with a new
identifier.  It is not, in general, observable by the biff server, which
sees only \texttt{CLAUDE\_PROJECT\_DIR} and the \texttt{initialize}
client name; a \texttt{SessionStart} hook captures it and correlates to
the server through the shared topmost-\texttt{claude} PID (biff-7ak).
The mechanism is code; the model takes as given only that a $SESSION$ is
a stable, non-recycled identity for the life of the Claude session,
surviving resume and distinct on fork.

\begin{zed}
  KIND ::= human | agent
\end{zed}

Every process is either a human operator or an AI agent.  $KIND$ is a
property of processes only.  A routing entry carries no role: one Claude
session hosts at most one routing entry, the entry of its agent.

%%--------------------------------------------------------------------
\subsection{The routing layer}
%%--------------------------------------------------------------------

Routing needs two further notions: a display alias and a choice of
routing discipline.

\begin{zed}
  [ALIAS]
\end{zed}

$ALIAS$ is the recyclable display name a teammate types --- the
\texttt{ttyN} of \texttt{@user:ttyN} (DES-035).  It is reserved
atomically and assigned lowest-free, so it is unique among the
concurrently-live entries but \emph{recycled} over time.  A message
payload is deliberately \emph{not} a type of this model: routing is about
\emph{where} a send goes, and, following DES-048, delivery correctness is
judged at the send rather than through a modelled recipient queue.

\begin{zed}
  ZBOOL ::= ztrue | zfalse
\end{zed}

\begin{zed}
  ROUTEMODE ::= rmIdentity | rmAlias
\end{zed}

$ROUTEMODE$ selects the routing discipline, fixed at initialization and
never changed --- a scoping parameter, as in $talk.tex$'s $routeMode$
(DES-048).  Under $rmIdentity$ --- the ratified design (biff-7ak,
Option~B) --- a send addressed to a display alias is resolved to the
entry that \emph{currently} holds that alias, \emph{at send time}, and
delivered to that entry's stable identity-keyed inbox.  Under $rmAlias$
--- the defect --- the send routes on a binding cached when the alias was
learned and never refreshed, so a resume or a recycle that has moved the
alias since strands the message (LOST) or lands it on the wrong live
entry (MISROUTED).

The routing identity of an entry is \emph{derived} from its session.

\begin{zed}
  RID == SESSION
\end{zed}

\begin{zed}
  derive == (\lambda s : SESSION @ s)
\end{zed}

$derive~s$ is the routing identity of the entry of session $s$.  In the
implementation it is a hex hash of the \texttt{session\_id}
(\texttt{blake2b(session\_id + role)}, with the role a single constant);
the hash is code and is not modelled.  What the model keeps is the whole
of what the hash is \emph{for}: derivation is a \emph{function}
(deterministic --- the same session always yields the same identity), it
is \emph{injective} (distinct sessions yield distinct identities), and it
depends on nothing but $s$, which survives \texttt{--resume} unchanged
(so the identity is stable across resume).  Modelling $derive$ as the
identity on $SESSION$ makes all three properties hold by construction;
because $derive$ is the identity function, its injectivity is trivial and
is not stated as an invariant.  This is the operator amendment to
biff-7ak: the entry carries a \emph{derived} identity, not a volatile
fresh hex, so it does not reproduce the bug.  An earlier version of this
model carried a second, companion role per session; the dual-session
companion was removed (PL-231) and with it the role component of the
routing identity.

%%--------------------------------------------------------------------
\section{State}
%%--------------------------------------------------------------------

The state captures the full hierarchy.  Each layer is modeled as a set
of entities plus partial functions mapping each entity to its attributes.
Invariants enforce referential integrity across layers.

\begin{schema}{BiffWorld}
  %% --- Organization layer ---
  orgs : \power NAME \\
  repos : NAME \rel NAME \\
  members : NAME \rel NAME \\
  %% --- Machine layer ---
  machines : \power HOST \\
  clones : HOST \rel (NAME \cross NAME) \\
  %% --- Session layer ---
  sessions : \power SESSION \\
  smach : SESSION \pfun HOST \\
  srepo : SESSION \pfun (NAME \cross NAME) \\
  sbranch : SESSION \pfun NAME \\
  soper : SESSION \pfun NAME \\
  %% --- Process layer ---
  processes : \power PROCESS \\
  psess : PROCESS \pfun SESSION \\
  pkind : PROCESS \pfun KIND \\
  pname : PROCESS \pfun NAME \\
  ppar : PROCESS \pfun PROCESS
  \where
  %%
  %% Organization scope: repos and members belong to existing orgs
  %%
  \dom repos \subseteq orgs \\
  \dom members \subseteq orgs \\
  %%
  %% Clone integrity: clones are on existing machines of existing repos
  %%
  \dom clones \subseteq machines \\
  \ran clones \subseteq repos \\
  %%
  %% Session attribute domains match the session set
  %%
  \dom smach = sessions \\
  \dom srepo = sessions \\
  \dom sbranch = sessions \\
  \dom soper \subseteq sessions \\
  %%
  %% Session references valid machine and clone
  %%
  \ran smach \subseteq machines \\
  \forall s : sessions @ (smach(s), srepo(s)) \in clones \\
  %%
  %% Operator is a member of the session's organization
  %%
  \forall s : \dom soper @
    (first(srepo(s)), soper(s)) \in members \\
  %%
  %% Process attribute domains match the process set
  %%
  \dom psess = processes \\
  \dom pkind = processes \\
  \dom pname = processes \\
  \dom ppar \subseteq processes \\
  %%
  %% Processes reference existing sessions; parents are processes
  %%
  \ran psess \subseteq sessions \\
  \ran ppar \subseteq processes \\
  %%
  %% Parent process is in the same session as the child
  %%
  \forall p : \dom ppar @
    psess(ppar(p)) = psess(p) \\
  %%
  %% Process tree is acyclic: no process is its own ancestor
  %%
  \forall p : processes @ (p, p) \notin ppar \plus
\end{schema}

%%--------------------------------------------------------------------
\section{Routing State}
%%--------------------------------------------------------------------

$RouteWorld$ extends $BiffWorld$ with the routing layer: the live
per-session entries, their recyclable display aliases, the
identity-keyed inboxes, and the discipline-selecting bookkeeping.  The
structural state is unchanged; the extension is where addressing and
delivery live.  Keeping it a distinct schema, included rather than
merged, lets the structural operations except $CloseSession$ remain framed
on $BiffWorld$ alone: they leave the routing components fixed, which the
animator supplies automatically.  Closing must also observe the routing
layer so that live entries are cleaned up before their session is removed.

\begin{schema}{RouteWorld}
  BiffWorld \\
  entries : \power SESSION \\
  ttyalias : SESSION \pinj ALIAS \\
  savedalias : SESSION \pfun ALIAS \\
  abind : ALIAS \pfun SESSION \\
  routeMode : ROUTEMODE \\
  lostF : ZBOOL \\
  misF : ZBOOL
  \where
  %%
  %% Every live entry belongs to a live session
  %%
  entries \subseteq sessions \\
  %%
  %% Every live entry has exactly one display alias, and --- the alias
  %% map being an injection --- no alias names two live entries at once
  %% (DES-035 uniqueness, lifted to entries)
  %%
  \dom ttyalias = entries \\
  %%
  %% Delivery soundness: under identity routing no message is ever
  %% stranded or misrouted (the biff-7ak guarantee)
  %%
  routeMode = rmIdentity \implies lostF = zfalse \\
  routeMode = rmIdentity \implies misF = zfalse
\end{schema}

$entries$ is the set of live routing endpoints, each identified by its
session; $entries \subseteq sessions$ ties every endpoint to a live
session.  $ttyalias$ carries each live entry to the display alias a
teammate types; declaring it an \emph{injection} ($\pinj$) is exactly
DES-035's global uniqueness lifted to entries --- a recycled alias
cannot name two live entries at once --- and $\dom ttyalias = entries$
says every live entry is displayable and every displayable name is live.
$savedalias$ remembers an entry's last alias across a suspension, so a
resume can reclaim it; unlike $ttyalias$ it is not cleared when the entry
goes, and it may name an entry that is no longer live.  $abind$ is the
cached ``alias~$\mapsto$~entry'' binding a route-on-alias send consults;
it is deliberately \emph{unconstrained} against $entries$, because its
staleness --- pointing at an entry that has since moved or gone --- is
the whole of the defect.  Message stores are not modelled: as in DES-048,
where $talk.tex$ proves routing soundness through the $unsoundEmit$ latch
rather than through recipient queues, delivery here is judged at the send
by comparing the coordinate a send resolves to against the entry that
holds the alias now.  $routeMode$ fixes the discipline for the run;
$lostF$ and $misF$ are history flags that latch when a send strands or
misroutes.  Because delivery under $rmIdentity$ resolves the alias at
send time (the entry that holds it \emph{now}), the identity-keyed inbox
it would file to is by construction the current holder's own; the latch,
not a stored queue, is what records whether that holds.  The last two
invariants are the delivery-soundness guarantee, stated so ProB decides
it as a state invariant: under $rmIdentity$ neither flag can ever latch.

Two facts the amendment turns on are here \emph{by construction} rather
than as invariants to be checked.  Routing identities are distinct by
construction because $derive$ is the identity on $SESSION$ and $entries$
is a set of distinct sessions: no two live entries share an identity, and
no re-derivation can collide them.  Stability across resume is likewise
structural: $ResumeSession$ re-mints the \emph{same} $s?$, so the identity
is literally unchanged.  What \emph{is} checked, and is the property a
recycled alias could otherwise break, is that $ttyalias$ is an injection
$SESSION \pinj ALIAS$ --- no display alias names two live entries at
once.

%%--------------------------------------------------------------------
\section{The Daemon Layer}
\label{sec:daemon}
%%--------------------------------------------------------------------

DES-095 puts the session records, the arrival subscription and the
unread state in one per-user daemon, \texttt{biffd}, and makes every
client a connection to it.  Three further notions are needed.

\begin{zed}
  [RUN, CONN]
\end{zed}

$RUN$ is the \texttt{main\_agent\_id} of a session: a new value for a
known session id is a new \emph{run} of that session (resume with a
fresh agent).  $CONN$ identifies one client connection to the daemon.
A connection is opened, becomes \emph{bound} to a session and run by a
valid \texttt{hello}, and may then \emph{subscribe}.

\begin{axdef}
  maxUnread : \nat
  \where
  maxUnread = 1
\end{axdef}

$maxUnread$ bounds the unread count so that ProB can enumerate it.  Mail
is durable, so the count of unread arrivals of a session is kept
\emph{outside} the daemon, with the relay: it survives the daemon.

\begin{schema}{DaemonSide}
  daemonUp : ZBOOL \\
  known : \power SESSION \\
  runOf : SESSION \pfun RUN \\
  opened : \power CONN \\
  hello : CONN \pfun (SESSION \cross RUN) \\
  subscribed : \power CONN \\
  shown : CONN \pfun (0 \upto maxUnread) \\
  durable : SESSION \fun (0 \upto maxUnread)
\end{schema}

\begin{schema}{DaemonSideInv}
  DaemonSide
  \where
  \dom runOf = known \\
  \dom hello \subseteq opened \\
  \forall c : \dom hello @ first(hello~c) \in known \\
  subscribed \subseteq \dom hello \\
  \forall c : subscribed @ second(hello~c) = runOf(first(hello~c)) \\
  \dom shown = subscribed \\
  \forall c : subscribed @ shown~c \leq durable(first(hello~c)) \\
  daemonUp = zfalse \implies opened = \emptyset
\end{schema}

$known$ is the set of session records the daemon holds, and $runOf$ the
current run of each.  A record is \emph{live} when some connection is
bound to it under its current run, and \emph{dormant} otherwise; there is
no third kind of record, and an ended session has none ($s \notin known$).
$opened$ is the set of open connections; $hello$ binds the
authenticated ones to a session and a run, so $\dom hello \subseteq
opened$ says that only an open connection can be bound, and a connection
in $opened \setminus \dom hello$ has sent no valid \texttt{hello}.
$subscribed$ is the set of bound connections with an arrival subscription,
and the quantified predicate over it is the replacement rule: a
subscription belongs to the \emph{current} run of its session, so the run
that a new \texttt{main\_agent\_id} has replaced holds none.
$shown~c$ is the unread summary that connection $c$ has been given, and
$durable~s$ the unread count that is kept for session $s$.  The pair is
the whole of invariant~4: $shown~c \leq durable(\ldots)$ says no
connection is ever told of more arrivals than exist (none duplicated),
and $Deliver$ and $Subscribe$ below say it is told of every one (none
lost).  When the daemon is not running there is no open connection, so no
record is live: a restart returns every record to dormant.

The predicates are stated in $DaemonSideInv$ and not in $DaemonSide$, which
carries only the declarations.  An operation written $\Delta DaemonSide$
includes the predicates of $DaemonSide$ on its after-state, and ProZ turns
them into guards: a transition that would break one is not enabled, and the
model checker cannot report it.  With the predicates in $DaemonSideInv$, which
$DaemonWorld$ includes and the reserved schema $Invariant$ installs, an
operation that breaks one is enabled, reaches the bad state, and is reported.

A connection holds no relay handle: $DaemonSide$ has no component of that
type, and the relay handle is a component of the daemon
(\texttt{nats-relay.tex}, $holdsHandle$).  Invariant~1 of DES-095 is
therefore a property of the shape of this schema, and is checked in the
relay specification, where the handle exists.

$DaemonWorld$ ties the daemon to the structural layer: a record exists
only for a session that exists.

\begin{schema}{DaemonWorld}
  RouteWorld \\
  DaemonSideInv
  \where
  known \subseteq sessions
\end{schema}

ProZ selects $DaemonWorld$ as the state through $Init$, but requires the
reserved schema name $Invariant$ to install its predicates as the
machine-checkable invariant.  Inclusion preserves all structural,
function, injection, delivery and daemon predicates without duplicating
them.

\begin{schema}{Invariant}
  DaemonWorld
\end{schema}

%%--------------------------------------------------------------------
\section{Initial State}
%%--------------------------------------------------------------------

The system starts empty: no organizations, machines, sessions,
processes, entries, aliases, or messages.

\begin{schema}{Init}
  DaemonWorld'
  \where
  orgs' = \emptyset \\
  repos' = \emptyset \\
  members' = \emptyset \\
  machines' = \emptyset \\
  clones' = \emptyset \\
  sessions' = \emptyset \\
  smach' = \emptyset \\
  srepo' = \emptyset \\
  sbranch' = \emptyset \\
  soper' = \emptyset \\
  processes' = \emptyset \\
  psess' = \emptyset \\
  pkind' = \emptyset \\
  pname' = \emptyset \\
  ppar' = \emptyset \\
  %%
  %% Routing layer starts empty; routeMode is left free, so ProB
  %% explores rmIdentity and rmAlias as two separate constant branches
  %%
  entries' = \emptyset \\
  ttyalias' = \emptyset \\
  savedalias' = \emptyset \\
  abind' = \emptyset \\
  lostF' = zfalse \\
  misF' = zfalse \\
  %%
  %% Daemon layer: running, no record, no connection, nothing unread
  %%
  daemonUp' = ztrue \\
  known' = \emptyset \\
  runOf' = \emptyset \\
  opened' = \emptyset \\
  hello' = \emptyset \\
  subscribed' = \emptyset \\
  shown' = \emptyset \\
  durable' = (\lambda s : SESSION @ 0)
\end{schema}

The daemon starts running with no session record, no open connection and
no unread arrival.

$routeMode'$ is left unconstrained: $Init$ admits both disciplines and no
operation reassigns it, so each discipline is a separate branch of the
state space.  Every latch starts $zfalse$, so the delivery-soundness
invariant holds at initialization under either discipline.

%%--------------------------------------------------------------------
\section{Organization Operations}
%%--------------------------------------------------------------------

%%--------------------------------------------------------------------
\subsection{AddOrg}
%%--------------------------------------------------------------------

Register a new organization.

\begin{schema}{AddOrg}
  \Delta BiffWorld \\
  org? : NAME
  \where
  org? \notin orgs \\
  %%
  orgs' = orgs \cup \{ org? \} \\
  %%
  %% Frame: everything else unchanged
  %%
  repos' = repos \\
  members' = members \\
  machines' = machines \\
  clones' = clones \\
  sessions' = sessions \\
  smach' = smach \\
  srepo' = srepo \\
  sbranch' = sbranch \\
  soper' = soper \\
  processes' = processes \\
  psess' = psess \\
  pkind' = pkind \\
  pname' = pname \\
  ppar' = ppar
\end{schema}

%%--------------------------------------------------------------------
\subsection{RemoveOrg}
%%--------------------------------------------------------------------

Remove an organization.  Precondition: no repos and no members remain.

\begin{schema}{RemoveOrg}
  \Delta BiffWorld \\
  org? : NAME
  \where
  org? \in orgs \\
  \{ org? \} \dres repos = \emptyset \\
  \{ org? \} \dres members = \emptyset \\
  %%
  orgs' = orgs \setminus \{ org? \} \\
  %%
  repos' = repos \\
  members' = members \\
  machines' = machines \\
  clones' = clones \\
  sessions' = sessions \\
  smach' = smach \\
  srepo' = srepo \\
  sbranch' = sbranch \\
  soper' = soper \\
  processes' = processes \\
  psess' = psess \\
  pkind' = pkind \\
  pname' = pname \\
  ppar' = ppar
\end{schema}

%%--------------------------------------------------------------------
\section{Repository Operations}
%%--------------------------------------------------------------------

%%--------------------------------------------------------------------
\subsection{AddRepo}
%%--------------------------------------------------------------------

Add a repository to an existing organization.

\begin{schema}{AddRepo}
  \Delta BiffWorld \\
  org? : NAME \\
  repo? : NAME
  \where
  org? \in orgs \\
  (org?, repo?) \notin repos \\
  %%
  repos' = repos \cup \{ (org?, repo?) \} \\
  %%
  orgs' = orgs \\
  members' = members \\
  machines' = machines \\
  clones' = clones \\
  sessions' = sessions \\
  smach' = smach \\
  srepo' = srepo \\
  sbranch' = sbranch \\
  soper' = soper \\
  processes' = processes \\
  psess' = psess \\
  pkind' = pkind \\
  pname' = pname \\
  ppar' = ppar
\end{schema}

%%--------------------------------------------------------------------
\subsection{RemoveRepo}
%%--------------------------------------------------------------------

Remove a repository.  Precondition: no clones of this repo exist on any
machine.

\begin{schema}{RemoveRepo}
  \Delta BiffWorld \\
  org? : NAME \\
  repo? : NAME
  \where
  (org?, repo?) \in repos \\
  (org?, repo?) \notin \ran clones \\
  %%
  repos' = repos \setminus \{ (org?, repo?) \} \\
  %%
  orgs' = orgs \\
  members' = members \\
  machines' = machines \\
  clones' = clones \\
  sessions' = sessions \\
  smach' = smach \\
  srepo' = srepo \\
  sbranch' = sbranch \\
  soper' = soper \\
  processes' = processes \\
  psess' = psess \\
  pkind' = pkind \\
  pname' = pname \\
  ppar' = ppar
\end{schema}

%%--------------------------------------------------------------------
\section{Machine Operations}
%%--------------------------------------------------------------------

%%--------------------------------------------------------------------
\subsection{AddMachine}
%%--------------------------------------------------------------------

Register a new machine.

\begin{schema}{AddMachine}
  \Delta BiffWorld \\
  m? : HOST
  \where
  m? \notin machines \\
  %%
  machines' = machines \cup \{ m? \} \\
  %%
  orgs' = orgs \\
  repos' = repos \\
  members' = members \\
  clones' = clones \\
  sessions' = sessions \\
  smach' = smach \\
  srepo' = srepo \\
  sbranch' = sbranch \\
  soper' = soper \\
  processes' = processes \\
  psess' = psess \\
  pkind' = pkind \\
  pname' = pname \\
  ppar' = ppar
\end{schema}

%%--------------------------------------------------------------------
\subsection{RemoveMachine}
%%--------------------------------------------------------------------

Remove a machine.  Precondition: no clones remain on it.

\begin{schema}{RemoveMachine}
  \Delta BiffWorld \\
  m? : HOST
  \where
  m? \in machines \\
  \{ m? \} \dres clones = \emptyset \\
  %%
  machines' = machines \setminus \{ m? \} \\
  %%
  orgs' = orgs \\
  repos' = repos \\
  members' = members \\
  clones' = clones \\
  sessions' = sessions \\
  smach' = smach \\
  srepo' = srepo \\
  sbranch' = sbranch \\
  soper' = soper \\
  processes' = processes \\
  psess' = psess \\
  pkind' = pkind \\
  pname' = pname \\
  ppar' = ppar
\end{schema}

%%--------------------------------------------------------------------
\section{Clone Operations}
%%--------------------------------------------------------------------

A clone represents a copy of a repository on a machine.  In practice
this is a \texttt{git clone} or a worktree checkout.

%%--------------------------------------------------------------------
\subsection{AddClone}
%%--------------------------------------------------------------------

Clone a repository onto a machine.

\begin{schema}{AddClone}
  \Delta BiffWorld \\
  m? : HOST \\
  org? : NAME \\
  repo? : NAME
  \where
  m? \in machines \\
  (org?, repo?) \in repos \\
  (m?, (org?, repo?)) \notin clones \\
  %%
  clones' = clones \cup \{ (m?, (org?, repo?)) \} \\
  %%
  orgs' = orgs \\
  repos' = repos \\
  members' = members \\
  machines' = machines \\
  sessions' = sessions \\
  smach' = smach \\
  srepo' = srepo \\
  sbranch' = sbranch \\
  soper' = soper \\
  processes' = processes \\
  psess' = psess \\
  pkind' = pkind \\
  pname' = pname \\
  ppar' = ppar
\end{schema}

%%--------------------------------------------------------------------
\subsection{RemoveClone}
%%--------------------------------------------------------------------

Remove a clone from a machine.  Precondition: no sessions are using
this clone.

\begin{schema}{RemoveClone}
  \Delta BiffWorld \\
  m? : HOST \\
  org? : NAME \\
  repo? : NAME
  \where
  (m?, (org?, repo?)) \in clones \\
  \lnot (\exists s : sessions @
    smach(s) = m? \land srepo(s) = (org?, repo?)) \\
  %%
  clones' = clones \setminus \{ (m?, (org?, repo?)) \} \\
  %%
  orgs' = orgs \\
  repos' = repos \\
  members' = members \\
  machines' = machines \\
  sessions' = sessions \\
  smach' = smach \\
  srepo' = srepo \\
  sbranch' = sbranch \\
  soper' = soper \\
  processes' = processes \\
  psess' = psess \\
  pkind' = pkind \\
  pname' = pname \\
  ppar' = ppar
\end{schema}

%%--------------------------------------------------------------------
\section{Member Operations}
%%--------------------------------------------------------------------

Members are user identities within an organization.  A member may
operate sessions and appear as a human process.

%%--------------------------------------------------------------------
\subsection{AddMember}
%%--------------------------------------------------------------------

Add a user to an organization.

\begin{schema}{AddMember}
  \Delta BiffWorld \\
  org? : NAME \\
  user? : NAME
  \where
  org? \in orgs \\
  (org?, user?) \notin members \\
  %%
  members' = members \cup \{ (org?, user?) \} \\
  %%
  orgs' = orgs \\
  repos' = repos \\
  machines' = machines \\
  clones' = clones \\
  sessions' = sessions \\
  smach' = smach \\
  srepo' = srepo \\
  sbranch' = sbranch \\
  soper' = soper \\
  processes' = processes \\
  psess' = psess \\
  pkind' = pkind \\
  pname' = pname \\
  ppar' = ppar
\end{schema}

%%--------------------------------------------------------------------
\subsection{RemoveMember}
%%--------------------------------------------------------------------

Remove a user from an organization.  Precondition: the user is not
operating any session on this org's repositories.

\begin{schema}{RemoveMember}
  \Delta BiffWorld \\
  org? : NAME \\
  user? : NAME
  \where
  (org?, user?) \in members \\
  \lnot (\exists s : \dom soper @
    first(srepo(s)) = org? \land soper(s) = user?) \\
  %%
  members' = members \setminus \{ (org?, user?) \} \\
  %%
  orgs' = orgs \\
  repos' = repos \\
  machines' = machines \\
  clones' = clones \\
  sessions' = sessions \\
  smach' = smach \\
  srepo' = srepo \\
  sbranch' = sbranch \\
  soper' = soper \\
  processes' = processes \\
  psess' = psess \\
  pkind' = pkind \\
  pname' = pname \\
  ppar' = ppar
\end{schema}

%%--------------------------------------------------------------------
\section{Session Operations}
%%--------------------------------------------------------------------

A session is a terminal on a machine, working on a specific branch of a
cloned repository.  It may have a human operator (typical Claude Code
usage) or be headless (CI, SDK, agent teams).

%%--------------------------------------------------------------------
\subsection{OpenSession}
%%--------------------------------------------------------------------

Open a new session with a human operator.  The operator must be a member
of the repository's organization.

\begin{schema}{OpenSession}
  \Delta BiffWorld \\
  s? : SESSION \\
  m? : HOST \\
  org? : NAME \\
  repo? : NAME \\
  branch? : NAME \\
  oper? : NAME
  \where
  s? \notin sessions \\
  (m?, (org?, repo?)) \in clones \\
  (org?, oper?) \in members \\
  %%
  sessions' = sessions \cup \{ s? \} \\
  smach' = smach \cup \{ s? \mapsto m? \} \\
  srepo' = srepo \cup \{ s? \mapsto (org?, repo?) \} \\
  sbranch' = sbranch \cup \{ s? \mapsto branch? \} \\
  soper' = soper \cup \{ s? \mapsto oper? \} \\
  %%
  orgs' = orgs \\
  repos' = repos \\
  members' = members \\
  machines' = machines \\
  clones' = clones \\
  processes' = processes \\
  psess' = psess \\
  pkind' = pkind \\
  pname' = pname \\
  ppar' = ppar
\end{schema}

%%--------------------------------------------------------------------
\subsection{OpenHeadlessSession}
%%--------------------------------------------------------------------

Open a session with no human operator.  This models CI runs, SDK
invocations, and autonomous agent teams.

\begin{schema}{OpenHeadlessSession}
  \Delta BiffWorld \\
  s? : SESSION \\
  m? : HOST \\
  org? : NAME \\
  repo? : NAME \\
  branch? : NAME
  \where
  s? \notin sessions \\
  (m?, (org?, repo?)) \in clones \\
  %%
  sessions' = sessions \cup \{ s? \} \\
  smach' = smach \cup \{ s? \mapsto m? \} \\
  srepo' = srepo \cup \{ s? \mapsto (org?, repo?) \} \\
  sbranch' = sbranch \cup \{ s? \mapsto branch? \} \\
  soper' = soper \\
  %%
  orgs' = orgs \\
  repos' = repos \\
  members' = members \\
  machines' = machines \\
  clones' = clones \\
  processes' = processes \\
  psess' = psess \\
  pkind' = pkind \\
  pname' = pname \\
  ppar' = ppar
\end{schema}

%%--------------------------------------------------------------------
\subsection{CloseSession}
%%--------------------------------------------------------------------

Close a session after all processes have been reaped and all live routing
entries have been suspended.  Live presence cleanup is automatic before
final close; this verification ordering adds no user-visible workflow
step.  $SuspendEntry$ releases each live alias while retaining routing
history; final close leaves that history unchanged.  A session that the
daemon still holds a record for ($s? \in known$) cannot be closed: the
record goes first, by $SessionEnd$ or $RetentionExpire$
(\S\ref{sec:daemonops}).

\begin{schema}{CloseSession}
  \Delta RouteWorld \\
  \Xi DaemonSide \\
  s? : SESSION
  \where
  s? \in sessions \\
  s? \notin \ran psess \\
  s? \notin entries \\
  s? \notin known \\
  %%
  sessions' = sessions \setminus \{ s? \} \\
  smach' = sessions' \dres smach \\
  srepo' = sessions' \dres srepo \\
  sbranch' = sessions' \dres sbranch \\
  soper' = sessions' \dres soper \\
  %%
  orgs' = orgs \\
  repos' = repos \\
  members' = members \\
  machines' = machines \\
  clones' = clones \\
  processes' = processes \\
  psess' = psess \\
  pkind' = pkind \\
  pname' = pname \\
  ppar' = ppar \\
  entries' = entries \\
  ttyalias' = ttyalias \\
  savedalias' = savedalias \\
  abind' = abind \\
  routeMode' = routeMode \\
  lostF' = lostF \\
  misF' = misF
\end{schema}

%%--------------------------------------------------------------------
\section{Process Operations}
%%--------------------------------------------------------------------

A process is any entity that can act within a session: a human operator
or an AI agent.  Processes form a tree.  Top-level processes have no
parent; subagents have a parent in the same session.

%%--------------------------------------------------------------------
\subsection{SpawnProcess}
%%--------------------------------------------------------------------

Spawn a new top-level process in a session.  This is used for the human
operator, the main agent, and agent team members.

\begin{schema}{SpawnProcess}
  \Delta BiffWorld \\
  p? : PROCESS \\
  s? : SESSION \\
  kind? : KIND \\
  name? : NAME
  \where
  p? \notin processes \\
  s? \in sessions \\
  %%
  processes' = processes \cup \{ p? \} \\
  psess' = psess \cup \{ p? \mapsto s? \} \\
  pkind' = pkind \cup \{ p? \mapsto kind? \} \\
  pname' = pname \cup \{ p? \mapsto name? \} \\
  ppar' = ppar \\
  %%
  orgs' = orgs \\
  repos' = repos \\
  members' = members \\
  machines' = machines \\
  clones' = clones \\
  sessions' = sessions \\
  smach' = smach \\
  srepo' = srepo \\
  sbranch' = sbranch \\
  soper' = soper
\end{schema}

%%--------------------------------------------------------------------
\subsection{SpawnChild}
%%--------------------------------------------------------------------

Spawn a child process (subagent) under an existing parent process.  The
child inherits the parent's session.

\begin{schema}{SpawnChild}
  \Delta BiffWorld \\
  p? : PROCESS \\
  parent? : PROCESS \\
  kind? : KIND \\
  name? : NAME
  \where
  p? \notin processes \\
  parent? \in processes \\
  %%
  processes' = processes \cup \{ p? \} \\
  psess' = psess \cup \{ p? \mapsto psess(parent?) \} \\
  pkind' = pkind \cup \{ p? \mapsto kind? \} \\
  pname' = pname \cup \{ p? \mapsto name? \} \\
  ppar' = ppar \cup \{ p? \mapsto parent? \} \\
  %%
  orgs' = orgs \\
  repos' = repos \\
  members' = members \\
  machines' = machines \\
  clones' = clones \\
  sessions' = sessions \\
  smach' = smach \\
  srepo' = srepo \\
  sbranch' = sbranch \\
  soper' = soper
\end{schema}

%%--------------------------------------------------------------------
\subsection{ReapProcess}
%%--------------------------------------------------------------------

Remove a process.  Precondition: no child processes remain.

\begin{schema}{ReapProcess}
  \Delta BiffWorld \\
  p? : PROCESS
  \where
  p? \in processes \\
  p? \notin \ran ppar \\
  %%
  processes' = processes \setminus \{ p? \} \\
  psess' = processes' \dres psess \\
  pkind' = processes' \dres pkind \\
  pname' = processes' \dres pname \\
  ppar' = processes' \dres ppar \\
  %%
  orgs' = orgs \\
  repos' = repos \\
  members' = members \\
  machines' = machines \\
  clones' = clones \\
  sessions' = sessions \\
  smach' = smach \\
  srepo' = srepo \\
  sbranch' = sbranch \\
  soper' = soper
\end{schema}

%%--------------------------------------------------------------------
\section{Routing and Delivery Operations}
%%--------------------------------------------------------------------

These operations change the routing layer and leave the structural state
alone.  Each is a $\Delta RouteWorld$ conjoined with $\Xi BiffWorld$:
the $\Xi$ pins all sixteen structural components unchanged in a single
line and brings the session set into scope for the preconditions, so
each schema below manages only the eight routing components.  The one
exception is $ForkSession$, which \emph{creates} a session and so frames
the structural components explicitly, in the manner of $OpenSession$.

Throughout, an entry's routing identity is $derive~s = s$, so
minting a fresh entry for a session it held before recovers the
\emph{same} identity automatically --- this is what makes resume lose
nothing.

%%--------------------------------------------------------------------
\subsection{RegisterEntry}
%%--------------------------------------------------------------------

Mint the routing entry of a live session, assigning it a fresh
display alias --- the lowest-free \texttt{ttyN}, modelled as any alias
not currently in use.  This is a first launch, or a \emph{recycle} when
the chosen alias is one a departed entry released.

\begin{schema}{RegisterEntry}
  \Delta RouteWorld \\
  \Xi BiffWorld \\
  s? : SESSION \\
  n? : ALIAS
  \where
  s? \in sessions \\
  s? \notin entries \\
  n? \notin \ran ttyalias \\
  %%
  entries' = entries \cup \{ s? \} \\
  ttyalias' = ttyalias \cup \{ s? \mapsto n? \} \\
  savedalias' = savedalias \oplus \{ s? \mapsto n? \} \\
  %%
  abind' = abind \\
  routeMode' = routeMode \\
  lostF' = lostF \\
  misF' = misF
\end{schema}

%%--------------------------------------------------------------------
\subsection{SuspendEntry}
%%--------------------------------------------------------------------

The hosting process exits cleanly.  The entry leaves $entries$ and its
alias is released for recycling, but its identity is durable: messages
addressed to $derive~s?$ still name it, and wait for the session to
resume.  $savedalias$ remembers the released alias.

\begin{schema}{SuspendEntry}
  \Delta RouteWorld \\
  \Xi BiffWorld \\
  s? : SESSION
  \where
  s? \in entries \\
  %%
  entries' = entries \setminus \{ s? \} \\
  ttyalias' = \{ s? \} \ndres ttyalias \\
  %%
  savedalias' = savedalias \\
  abind' = abind \\
  routeMode' = routeMode \\
  lostF' = lostF \\
  misF' = misF
\end{schema}

%%--------------------------------------------------------------------
\subsection{ResumeSession}
%%--------------------------------------------------------------------

A suspended entry comes back.  Because $derive~s?$ depends only on
the session, the re-minted entry has the \emph{same}
routing identity --- so every message addressed to it before the
suspension is still addressed to it after, and nothing is stranded.  The
display alias is reclaimed when it is still free
($n? = savedalias~s?$), and otherwise a fresh free alias is taken.
Resume is per session: a full session resume fires this once, and the
identity persists in every case.

\begin{schema}{ResumeSession}
  \Delta RouteWorld \\
  \Xi BiffWorld \\
  s? : SESSION \\
  n? : ALIAS
  \where
  s? \in sessions \\
  s? \notin entries \\
  s? \in \dom savedalias \\
  n? \notin \ran ttyalias \\
  ( savedalias~s? \notin \ran ttyalias \implies
    n? = savedalias~s? ) \\
  %%
  entries' = entries \cup \{ s? \} \\
  ttyalias' = ttyalias \cup \{ s? \mapsto n? \} \\
  savedalias' = savedalias \oplus \{ s? \mapsto n? \} \\
  %%
  abind' = abind \\
  routeMode' = routeMode \\
  lostF' = lostF \\
  misF' = misF
\end{schema}

%%--------------------------------------------------------------------
\subsection{ForkSession}
%%--------------------------------------------------------------------

A fork mints a \emph{new} session identity, distinct from its parent, in
the same clone, and gives it a fresh agent entry with a fresh alias.
Because the session is new, $derive~s?$ is a routing identity no
entry has ever held --- a fork inherits none of the parent's mail.  This
is the one routing operation that changes the structural state, so it
frames the structural components in full.

\begin{schema}{ForkSession}
  \Delta RouteWorld \\
  sp? : SESSION \\
  s? : SESSION \\
  m? : HOST \\
  org? : NAME \\
  repo? : NAME \\
  branch? : NAME \\
  na? : ALIAS
  \where
  sp? \in sessions \\
  s? \notin sessions \\
  s? \neq sp? \\
  smach~sp? = m? \\
  srepo~sp? = (org?, repo?) \\
  na? \notin \ran ttyalias \\
  %%
  sessions' = sessions \cup \{ s? \} \\
  smach' = smach \cup \{ s? \mapsto m? \} \\
  srepo' = srepo \cup \{ s? \mapsto (org?, repo?) \} \\
  sbranch' = sbranch \cup \{ s? \mapsto branch? \} \\
  soper' = soper \\
  %%
  entries' = entries \cup \{ s? \} \\
  ttyalias' = ttyalias \cup \{ s? \mapsto na? \} \\
  savedalias' = savedalias \oplus \{ s? \mapsto na? \} \\
  abind' = abind \\
  routeMode' = routeMode \\
  lostF' = lostF \\
  misF' = misF \\
  %%
  orgs' = orgs \\
  repos' = repos \\
  members' = members \\
  machines' = machines \\
  clones' = clones \\
  processes' = processes \\
  psess' = psess \\
  pkind' = pkind \\
  pname' = pname \\
  ppar' = ppar
\end{schema}

%%--------------------------------------------------------------------
\subsection{Learn}
%%--------------------------------------------------------------------

A teammate resolves a display alias to an entry --- the \texttt{who} that
populates the binding a route-on-alias send will later consult.  It
records the \emph{current} holder; nothing refreshes it afterwards, so a
subsequent resume or recycle leaves it stale.

\begin{schema}{Learn}
  \Delta RouteWorld \\
  \Xi BiffWorld \\
  n? : ALIAS
  \where
  n? \in \ran ttyalias \\
  %%
  abind' = abind \oplus \{ n? \mapsto (ttyalias\inv)~n? \} \\
  %%
  entries' = entries \\
  ttyalias' = ttyalias \\
  savedalias' = savedalias \\
  routeMode' = routeMode \\
  lostF' = lostF \\
  misF' = misF
\end{schema}

%%--------------------------------------------------------------------
\subsection{DeliverIdentity}
%%--------------------------------------------------------------------

The ratified discipline.  A send addressed to alias $to?$ resolves it to
the entry that holds $to?$ \emph{now} --- $(ttyalias\inv)~to?$, well
defined because $to? \in \ran ttyalias$ --- and delivers to that entry's
identity-keyed inbox.  The predicate $(ttyalias\inv)~to? \in entries$
records that the resolution, re-evaluated at send time, always names a
\emph{live} current holder (it is a theorem of the $RouteWorld$
invariant $\dom ttyalias = entries$, not an added assumption); so the
send reaches the current holder and neither latch can fire.  No payload
is modelled --- routing is about \emph{where} a send goes.

\begin{schema}{DeliverIdentity}
  \Delta RouteWorld \\
  \Xi BiffWorld \\
  to? : ALIAS
  \where
  routeMode = rmIdentity \\
  to? \in \ran ttyalias \\
  (ttyalias\inv)~to? \in entries \\
  %%
  entries' = entries \\
  ttyalias' = ttyalias \\
  savedalias' = savedalias \\
  abind' = abind \\
  routeMode' = routeMode \\
  lostF' = lostF \\
  misF' = misF
\end{schema}

%%--------------------------------------------------------------------
\subsection{DeliverAliasHit}
%%--------------------------------------------------------------------

The defective discipline, in the case where the cached binding happens
still to be right.  A route-on-alias send routes to $abind~to?$; when
that entry is live and still holds $to?$, the send reaches the intended
holder, correct by luck.  It is the first of three cases that partition
the route-on-alias send.

\begin{schema}{DeliverAliasHit}
  \Delta RouteWorld \\
  \Xi BiffWorld \\
  to? : ALIAS
  \where
  routeMode = rmAlias \\
  to? \in \dom abind \\
  abind~to? \in entries \\
  ttyalias~(abind~to?) = to? \\
  %%
  entries' = entries \\
  ttyalias' = ttyalias \\
  savedalias' = savedalias \\
  abind' = abind \\
  routeMode' = routeMode \\
  lostF' = lostF \\
  misF' = misF
\end{schema}

%%--------------------------------------------------------------------
\subsection{DeliverAliasLost}
%%--------------------------------------------------------------------

The LOST case.  The cached binding names an entry that is no longer live
--- the session was resumed under a new coordinate, or has gone --- so a
route-on-alias send routed by that binding reaches no live entry and
strands; $lostF$ latches.  This is reachable only under $rmAlias$:
identity routing would re-resolve $to?$ and either reach the current
holder or, if none, decline to send.

\begin{schema}{DeliverAliasLost}
  \Delta RouteWorld \\
  \Xi BiffWorld \\
  to? : ALIAS
  \where
  routeMode = rmAlias \\
  to? \in \dom abind \\
  abind~to? \notin entries \\
  %%
  lostF' = ztrue \\
  %%
  entries' = entries \\
  ttyalias' = ttyalias \\
  savedalias' = savedalias \\
  abind' = abind \\
  routeMode' = routeMode \\
  misF' = misF
\end{schema}

%%--------------------------------------------------------------------
\subsection{DeliverAliasMis}
%%--------------------------------------------------------------------

The MISROUTED case.  The cached binding names an entry that is still live
but has since moved off the alias --- the alias was recycled to a
different session, or reassigned on resume --- so a route-on-alias send
lands on a live but \emph{wrong} entry.  $ttyalias~(abind~to?) \neq to?$
is well defined because $abind~to? \in entries = \dom ttyalias$, and says
exactly that the cached target no longer holds the alias it was addressed
by.  The send reaches the wrong entry; $misF$ latches.  Reachable only
under $rmAlias$.

\begin{schema}{DeliverAliasMis}
  \Delta RouteWorld \\
  \Xi BiffWorld \\
  to? : ALIAS
  \where
  routeMode = rmAlias \\
  to? \in \dom abind \\
  abind~to? \in entries \\
  ttyalias~(abind~to?) \neq to? \\
  %%
  misF' = ztrue \\
  %%
  entries' = entries \\
  ttyalias' = ttyalias \\
  savedalias' = savedalias \\
  abind' = abind \\
  routeMode' = routeMode \\
  lostF' = lostF
\end{schema}

%%--------------------------------------------------------------------
\section{Daemon Operations}
\label{sec:daemonops}
%%--------------------------------------------------------------------

Every operation below is a $\Delta DaemonSide$ conjoined with $\Xi
RouteWorld$: it leaves the structural and routing layers alone and frames
each of the eight daemon components explicitly.  The frames are the
content.  An operation that does not mention $known$, $runOf$ or
$durable$ in its effect leaves them as they were, and that is how
``detach is not session end'' and ``a restart loses no mail'' are stated.

%%--------------------------------------------------------------------
\subsection{Daemon lifecycle}
%%--------------------------------------------------------------------

\begin{schema}{DaemonExit}
  \Delta DaemonSide \\
  \Xi RouteWorld
  \where
  daemonUp = ztrue \\
  daemonUp' = zfalse \\
  opened' = \emptyset \\
  hello' = \emptyset \\
  subscribed' = \emptyset \\
  shown' = \emptyset \\
  %%
  known' = known \\
  runOf' = runOf \\
  durable' = durable
\end{schema}

The process ends, by a clean stop or a crash.  Every connection goes, so
every record is dormant, and nothing else changes: the records, their
runs and the unread counts are kept ($known'$, $runOf'$, $durable'$).  A
restart therefore finds every session where it was, with no client
attached, and loses no arrival.

\begin{schema}{DaemonStart}
  \Delta DaemonSide \\
  \Xi RouteWorld
  \where
  daemonUp = zfalse \\
  daemonUp' = ztrue \\
  %%
  known' = known \\
  runOf' = runOf \\
  opened' = opened \\
  hello' = hello \\
  subscribed' = subscribed \\
  shown' = shown \\
  durable' = durable
\end{schema}

The service manager, or a client's autostart nudge, starts the daemon.
It starts with no connection, because $daemonUp = zfalse \implies opened =
\emptyset$ held before it.

\begin{schema}{Arrive}
  \Delta DaemonSide \\
  \Xi RouteWorld \\
  s? : SESSION
  \where
  durable~s? < maxUnread \\
  durable' = durable \oplus \{ s? \mapsto durable~s? + 1 \} \\
  %%
  daemonUp' = daemonUp \\
  known' = known \\
  runOf' = runOf \\
  opened' = opened \\
  hello' = hello \\
  subscribed' = subscribed \\
  shown' = shown
\end{schema}

A message reaches the identity-keyed inbox of $s?$.  It is enabled whether
or not the daemon is running and whether or not the session has a record:
the relay keeps the mail.  Only $durable$ moves; telling a client is the
separate step $Deliver$.

%%--------------------------------------------------------------------
\subsection{Connections and the hello}
%%--------------------------------------------------------------------

\begin{schema}{Connect}
  \Delta DaemonSide \\
  \Xi RouteWorld \\
  c? : CONN
  \where
  daemonUp = ztrue \\
  c? \notin opened \\
  opened' = opened \cup \{ c? \} \\
  %%
  daemonUp' = daemonUp \\
  known' = known \\
  runOf' = runOf \\
  hello' = hello \\
  subscribed' = subscribed \\
  shown' = shown \\
  durable' = durable
\end{schema}

A client opens a connection.  It is not yet bound to any session.

\begin{schema}{RequestBeforeHello}
  \Delta DaemonSide \\
  \Xi RouteWorld \\
  c? : CONN
  \where
  c? \in opened \\
  c? \notin \dom hello \\
  opened' = opened \setminus \{ c? \} \\
  %%
  daemonUp' = daemonUp \\
  known' = known \\
  runOf' = runOf \\
  hello' = hello \\
  subscribed' = subscribed \\
  shown' = shown \\
  durable' = durable
\end{schema}

A request, or a \texttt{hello} with a wrong token, on a connection that has
no valid \texttt{hello}.  The daemon closes the connection and changes
nothing else: every component of the session layer other than the
rejected connection itself is as it was.  This is invariant~2 of DES-095.

\begin{schema}{HelloCreate}
  \Delta DaemonSide \\
  \Xi RouteWorld \\
  c? : CONN \\
  s? : SESSION \\
  r? : RUN
  \where
  c? \in opened \\
  c? \notin \dom hello \\
  s? \in sessions \\
  s? \notin known \\
  known' = known \cup \{ s? \} \\
  runOf' = runOf \cup \{ s? \mapsto r? \} \\
  hello' = hello \cup \{ c? \mapsto (s?, r?) \} \\
  %%
  daemonUp' = daemonUp \\
  opened' = opened \\
  subscribed' = subscribed \\
  shown' = shown \\
  durable' = durable
\end{schema}

The first \texttt{hello} for a session id creates its record, bound to the
run it names.

\begin{schema}{HelloAttach}
  \Delta DaemonSide \\
  \Xi RouteWorld \\
  c? : CONN \\
  s? : SESSION \\
  r? : RUN
  \where
  c? \in opened \\
  c? \notin \dom hello \\
  s? \in known \\
  runOf~s? = r? \\
  hello' = hello \cup \{ c? \mapsto (s?, r?) \} \\
  %%
  daemonUp' = daemonUp \\
  known' = known \\
  runOf' = runOf \\
  opened' = opened \\
  subscribed' = subscribed \\
  shown' = shown \\
  durable' = durable
\end{schema}

A client for a known session, in the run the record holds, attaches to the
existing record.  $known'$ and $runOf'$ are unchanged, so the daemon never
holds two records for one session id; $subscribed'$ and $shown'$ are
unchanged, so a second client attaching does not disturb the subscription
of the first.

\begin{schema}{HelloNewRun}
  \Delta DaemonSide \\
  \Xi RouteWorld \\
  c? : CONN \\
  s? : SESSION \\
  r? : RUN
  \where
  c? \in opened \\
  c? \notin \dom hello \\
  s? \in known \\
  runOf~s? \neq r? \\
  runOf' = runOf \oplus \{ s? \mapsto r? \} \\
  hello' = hello \cup \{ c? \mapsto (s?, r?) \} \\
  subscribed' = subscribed \setminus
    \{ c : subscribed | first(hello~c) = s? \} \\
  shown' = \{ c : subscribed | first(hello~c) = s? \} \ndres shown \\
  %%
  daemonUp' = daemonUp \\
  known' = known \\
  opened' = opened \\
  durable' = durable
\end{schema}

A new \texttt{main\_agent\_id} for a known session is a new run.  The
record is kept ($known'$, $durable'$) and takes the new run.  The previous
run's subscription and its pending deliveries are released: every
subscribed connection of the session was subscribed under the old run, so
$subscribed'$ and $shown'$ drop them.  The old run's connection stays open
and bound, so that the daemon can answer it with a named error
($RequestReplacedRun$), but it holds no subscription.

\begin{schema}{RequestReplacedRun}
  \Xi DaemonSide \\
  \Xi RouteWorld \\
  c? : CONN
  \where
  c? \in \dom hello \\
  second(hello~c?) \neq runOf(first(hello~c?))
\end{schema}

A call from a replaced run is refused with a named error and changes
nothing.

%%--------------------------------------------------------------------
\subsection{Subscription and unread state}
%%--------------------------------------------------------------------

\begin{schema}{Subscribe}
  \Delta DaemonSide \\
  \Xi RouteWorld \\
  c? : CONN
  \where
  c? \in \dom hello \\
  c? \notin subscribed \\
  second(hello~c?) = runOf(first(hello~c?)) \\
  subscribed' = subscribed \cup \{ c? \} \\
  shown' = shown \cup \{ c? \mapsto durable(first(hello~c?)) \} \\
  %%
  daemonUp' = daemonUp \\
  known' = known \\
  runOf' = runOf \\
  opened' = opened \\
  hello' = hello \\
  durable' = durable
\end{schema}

The current run's connection subscribes.  The daemon first sends the
session's current unread summary, so $shown'~c? = durable(\ldots)$: a
client that reattaches after a restart is told exactly the durable unread
count, and $durable'$ is unchanged, so the restart neither lost an arrival
nor counted one twice.  This is invariant~4 of DES-095.

\begin{schema}{Deliver}
  \Delta DaemonSide \\
  \Xi RouteWorld \\
  c? : CONN
  \where
  c? \in subscribed \\
  shown~c? < durable(first(hello~c?)) \\
  shown' = shown \oplus \{ c? \mapsto shown~c? + 1 \} \\
  %%
  daemonUp' = daemonUp \\
  known' = known \\
  runOf' = runOf \\
  opened' = opened \\
  hello' = hello \\
  subscribed' = subscribed \\
  durable' = durable
\end{schema}

One pending arrival is told to a subscribed connection.  It raises the
summary by one and no more, so a connection is told of each arrival
once.  A pending arrival is one with $shown~c < durable(\ldots)$.

\begin{schema}{ReadUnread}
  \Delta DaemonSide \\
  \Xi RouteWorld \\
  c? : CONN
  \where
  c? \in \dom hello \\
  second(hello~c?) = runOf(first(hello~c?)) \\
  durable' = durable \oplus \{ first(hello~c?) \mapsto 0 \} \\
  shown' =
    \{ c : subscribed | first(hello~c) = first(hello~c?) @ c \mapsto 0 \}
    \cup
    \{ c : subscribed | first(hello~c) \neq first(hello~c?) @
        c \mapsto shown~c \} \\
  %%
  daemonUp' = daemonUp \\
  known' = known \\
  runOf' = runOf \\
  opened' = opened \\
  hello' = hello \\
  subscribed' = subscribed
\end{schema}

Reading consumes the session's unread mail --- read-once, one writer ---
and every subscribed connection of that session is told the new, empty
summary.  The reader need not be subscribed: a one-shot read is a request
from a bound connection of the current run, as a call from a replaced run
is refused ($RequestReplacedRun$).

%%--------------------------------------------------------------------
\subsection{Detach and the end of a session}
%%--------------------------------------------------------------------

\begin{schema}{Detach}
  \Delta DaemonSide \\
  \Xi RouteWorld \\
  c? : CONN
  \where
  c? \in opened \\
  opened' = opened \setminus \{ c? \} \\
  hello' = \{ c? \} \ndres hello \\
  subscribed' = subscribed \setminus \{ c? \} \\
  shown' = \{ c? \} \ndres shown \\
  %%
  daemonUp' = daemonUp \\
  known' = known \\
  runOf' = runOf \\
  durable' = durable
\end{schema}

A dropped connection.  $known'$, $runOf'$ and $durable'$ are unchanged:
detaching ends no session, loses no mail and forgets no run.  The record
simply becomes dormant when its last connection goes.

\begin{schema}{SessionEnd}
  \Delta DaemonSide \\
  \Xi RouteWorld \\
  c? : CONN \\
  s? : SESSION
  \where
  s? \in known \\
  c? \in \dom hello \\
  first(hello~c?) = s? \\
  second(hello~c?) = runOf~s? \\
  known' = known \setminus \{ s? \} \\
  runOf' = \{ s? \} \ndres runOf \\
  opened' = opened \setminus
    \{ c : \dom hello | first(hello~c) = s? \} \\
  hello' = \{ c : \dom hello | first(hello~c) = s? \} \ndres hello \\
  subscribed' = subscribed \setminus
    \{ c : \dom hello | first(hello~c) = s? \} \\
  shown' = \{ c : \dom hello | first(hello~c) = s? \} \ndres shown \\
  %%
  daemonUp' = daemonUp \\
  durable' = durable
\end{schema}

The \texttt{SessionEnd} hook, which arrives over the connection $c?$,
ends the session only if $c?$ is bound to the session's \emph{current} run:
a stale hook from a replaced run, whose connection is still open, is refused
and leaves the session and the current connection intact.  The record goes
and so do its connections.  The unread count of the identity-keyed inbox stays.  This is
one of the two operations that remove a record.

\begin{schema}{RetentionExpire}
  \Delta DaemonSide \\
  \Xi RouteWorld \\
  s? : SESSION
  \where
  s? \in known \\
  \lnot ( \exists c : \dom hello @ first(hello~c) = s? ) \\
  known' = known \setminus \{ s? \} \\
  runOf' = \{ s? \} \ndres runOf \\
  %%
  daemonUp' = daemonUp \\
  opened' = opened \\
  hello' = hello \\
  subscribed' = subscribed \\
  shown' = shown \\
  durable' = durable
\end{schema}

A crashed session, with no client for the whole retention period, expires.
Only a dormant record can expire.  This is the other operation that
removes a record; with $SessionEnd$ it is the only way a session ends, so
detach, daemon exit and a new run all keep the record.

%%--------------------------------------------------------------------
\section{Delivery Properties and Verification}
\label{sec:delivery}
%%--------------------------------------------------------------------

Three properties are owed to biff-7ak.  Each is stated so that ProB
decides it, the first as a state invariant and the bugs as reachability
goals, exactly as DES-048 verifies $talk.tex$.

\begin{enumerate}
\item \textbf{Identity delivery.}  A send to a live alias under
  $rmIdentity$ resolves to the entry that holds the alias at send time,
  $(ttyalias\inv)~to?$, and to no other coordinate; and that entry is
  live ($DeliverIdentity$ asserts $(ttyalias\inv)~to? \in entries$).
  Delivery is a function of the identity, never of the alias string.

\item \textbf{Resume loses nothing.}  $ResumeSession$ re-mints the entry
  at the same routing identity $derive~s? = s?$, so every send
  addressed to $s?$ before the suspension resolves to it after the
  resume.  No send is stranded by a resume --- $lostF$ cannot latch under
  identity routing.

\item \textbf{Recycled alias never delivers to a stale entry.}  If
  $ttyalias$ carried alias $n$ to entry $e_1$, $e_1$ released it, and a
  later entry $e_2$ claimed $n$ (recycle), then under $rmIdentity$ a send
  to $n$ resolves $(ttyalias\inv)~n = e_2$ at send time and reaches $e_2$,
  never the stale $e_1$.
\end{enumerate}

All three are consequences of one state invariant of $RouteWorld$,
$$routeMode = rmIdentity \implies (lostF = zfalse \land misF = zfalse),$$
together with the definition of $DeliverIdentity$: under identity routing
every send resolves to the current, live holder, so none strands
(property~2 and the liveness half of property~3) and none misroutes
(property~1 and the correctness half of property~3).  The two defects are
the two ways a flag latches under $rmAlias$: $DeliverAliasLost$ is LOST
(property~2 violated --- a send to a resumed or departed session strands)
and $DeliverAliasMis$ is MISROUTED (properties~1 and~3 violated --- a send
reaches a live but wrong entry).  That no two live entries ever share a
routing identity --- in particular no old and new occupant of a recycled
alias --- follows by construction, identities being distinct sessions; and that no display alias
names two live entries at once is the injection
$ttyalias : SESSION \pinj ALIAS$, checked by ProB in every
reachable state.

\subsection{Model-checking scope}

The structural layer explodes at a uniform set size (two names alone put
the organization, repository, member and branch combinatorics past a
hundred and forty thousand states), so, as $talk.tex$ does, the routing
properties are checked at an explicit mixed scope that collapses the
structural carriers and keeps the routing carriers just large enough to
exhibit the defect.  The primary run is exhaustive at one session and two
aliases: one session carries one routing entry, and one alias suffices
to recycle while the other is the one the entry moves to on resume, so a
single session exercises both a stranded send and a misrouted one --- the
recycle happens across a suspend and a re-register of that one entry,
which resumes onto a fresh alias.

\begin{verbatim}
probcli docs/specs/session-model.tex -model_check -strict \
  -card NAME 1 -card HOST 1 -card PROCESS 1 \
  -card SESSION 1 -card ALIAS 2 \
  -p DEFAULT_SETSIZE 1 -p MAXINT 4 \
  -p MAX_INITIALISATIONS 100 -p MAX_OPERATIONS 20000 \
  -p TIME_OUT 60000 -coverage
\end{verbatim}

With the explicit $Invariant$ and cleanup-before-close guard, this visits
all $1{,}063$ reachable states ($5{,}372$ transitions) with
$invariant\_violated = 0$ and no deadlock, so every reachable state holds
the $RouteWorld$ invariant --- in particular $routeMode = rmIdentity
\implies lostF = zfalse \land misF = zfalse$, and the injectivity of
$ttyalias$ over live entries ($SESSION \pinj ALIAS$).  $ForkSession$ and $SpawnChild$ are uncovered:
they require a second $SESSION$ and $PROCESS$, respectively.  No larger
scope or fork/child coverage is claimed.  The uniform-size-one gate also
passes constraint-based deadlock checking and exhaustively explores
$193$ states ($788$ transitions).

Historical runs reported $15{,}247$ mixed-scope states and $966$
uniform-size-one states without an explicit $Invariant$.  Their translated
global invariant checked only underlying types, not the full $RouteWorld$
predicates, so those results did not establish the claimed invariants.
Installing the full invariant exposed $CloseSession$ leaving a live entry
whose owning session had been removed; cleanup-before-close repairs that
transition without weakening any invariant.

\medskip\noindent\textbf{(D1) Soundness holds --- as a state invariant.}
Under $rmIdentity$ neither latch is ever set: the exhaustive run above
finds no counterexample to the full $RouteWorld$ invariant across all
$1{,}063$ states at the stated bounds.

\medskip\noindent\textbf{(D2) LOST is reachable under alias routing.}
As a reachability goal,

\begin{verbatim}
probcli ... -card SESSION 1 -card ALIAS 2 -goal "lostF=ztrue & routeMode=rmAlias"
\end{verbatim}

ProB finds the goal.  The trace: $OpenHeadlessSession(s_1)$;
$RegisterEntry(s_1, n_2)$ mints the entry on $n_2$;
$Learn(n_2)$ caches $abind~n_2 = s_1$; $SuspendEntry(s_1)$
takes the entry out of $entries$; a route-on-alias send to $n_2$ then
matches $DeliverAliasLost$ ($abind~n_2 \notin entries$) and strands.  This
is the biff-7ak LOST bug: mail sent to the address a teammate holds lands
where no live session now drains.

\medskip\noindent\textbf{(D3) MISROUTED is reachable under alias
routing.}  As a reachability goal,

\begin{verbatim}
probcli ... -card SESSION 1 -card ALIAS 2 -goal "misF=ztrue & routeMode=rmAlias"
\end{verbatim}

ProB finds the goal.  The trace: $OpenHeadlessSession(s_1)$;
$RegisterEntry(s_1, n_2)$; $Learn(n_2)$ caches $abind~n_2 =
s_1$; $SuspendEntry(s_1)$ frees $n_2$;
$RegisterEntry(s_1, n_1)$ brings the entry back on a
\emph{different} alias (a resume that did not reclaim its old name).  A
route-on-alias send to $n_2$ then matches $DeliverAliasMis$
($abind~n_2 = s_1 \in entries$ but $ttyalias~s_1 = n_1
\neq n_2$) and lands on the live but wrong entry.  This is the biff-7ak
MISROUTED bug: a stale \texttt{ttyN} carries mail to a session that is
not its addressee.

\medskip\noindent\textbf{(D4) Both bugs are unreachable under identity
routing.}  As a reachability goal,

\begin{verbatim}
probcli ... -card SESSION 1 -card ALIAS 2 \
  -goal "(lostF=ztrue or misF=ztrue) & routeMode=rmIdentity"
\end{verbatim}

ProB reports no such state after visiting all $1{,}063$: under identity
routing neither latch can be set.  (D1) and (D4) are the same fact seen
twice --- as an invariant that holds and as a bad state that is
unreachable --- which is the precise sense in which routing on the
derived per-session identity makes the strand and the misroute
\emph{impossible} rather than merely unobserved.

%%--------------------------------------------------------------------
\section{Discussion}
%%--------------------------------------------------------------------

\subsection{What This Model Captures}

The hierarchy from organization to process is fully represented.  Key
properties enforced by the invariants:

\begin{itemize}
  \item A session can only exist on a machine that has a clone of the
    repository it references.
  \item An operator must be a member of the repository's organization.
  \item A headless session (no operator) is a first-class citizen, not
    an edge case.
  \item Subagent processes are children of their spawning agent, always
    in the same session.
  \item The process tree is acyclic.
  \item Removal at any level requires dependents to be removed first,
    preventing dangling references.
\end{itemize}

\subsection{Scenarios}

\noindent\textbf{Pair programming (typical Claude Code):}

\begin{verbatim}
  OpenSession(s1, macbook, punt-labs, biff, main, kai)
  SpawnProcess(p1, s1, human, kai)
  SpawnProcess(p2, s1, agent, claude)
  SpawnChild(p3, p2, agent, explore)   -- subagent
  SpawnChild(p4, p2, agent, explore)   -- another subagent
\end{verbatim}

\noindent\textbf{Headless CI:}

\begin{verbatim}
  OpenHeadlessSession(s2, ci-runner, punt-labs, biff, main)
  SpawnProcess(p5, s2, agent, claude)
  SpawnChild(p6, p5, agent, test-runner)
\end{verbatim}

\noindent\textbf{Agent team (no human):}

\begin{verbatim}
  OpenHeadlessSession(s3, workstation, punt-labs, biff, feature-x)
  SpawnProcess(p7, s3, agent, architect)
  SpawnProcess(p8, s3, agent, implementer)
  SpawnProcess(p9, s3, agent, reviewer)
\end{verbatim}

\noindent\textbf{User across machines and branches:}

\begin{verbatim}
  OpenSession(s4, macbook, punt-labs, biff, main, kai)
  OpenSession(s5, macbook, punt-labs, biff, feature-x, kai)
  OpenSession(s6, server, punt-labs, quarry, main, kai)
\end{verbatim}

\subsection{What This Model Does Not Cover}

This specification covers the structural containment hierarchy and the
routing layer --- the identity-keyed addressing and delivery of
biff-7ak.  It deliberately omits:

\begin{itemize}
  \item \textbf{Presence} --- idle times, plans, active status, and the
    display pipeline (\texttt{who}, \texttt{finger}, \texttt{ps}).
  \item \textbf{Permissions} --- who can message whom, who can steer
    whose session, \texttt{mesg y/n}.
  \item \textbf{Process lifecycle details} --- heartbeats, TTLs, orphan
    detection, crash recovery.
  \item \textbf{Implementation of derivation and correlation} --- the
    \texttt{blake2b} hex hash that realizes $derive$, the
    \texttt{SessionStart}-hook-to-server correlation over the shared
    \texttt{claude} PID, the \texttt{(pid, start-time)} recycle guard,
    and the KV storage of the resume-reclaim mapping.  These are code;
    the model keeps only the properties they must satisfy --- a
    deterministic, injective, resume-stable routing identity per
    session, and an alias that a send resolves at send time.
\end{itemize}

These remaining omissions are natural extensions: presence is a derived
view over the same state, and permissions a gate on the delivery
operations.

\end{document}
